%% ================================================================
%% SFST v30 — Neuer Supplement-Anhang: B1 Numerische Untersuchung
%% Titel: Appendix S: Numerical investigation of the second-order
%%         electromagnetic correction
%%
%% Einzufügen nach Appendix R (B3+B4).
%%
%% Tier-Stufen:
%%   Computation of V_eff^(2)(R0) / V_eff^(0)(R0): TIER 1
%%   Interpretation as α²-scale confirmation: TIER 1
%%   Non-derivability of 1/√8 from simple V_eff: TIER 1
%%   What would close B1: identified but not executed
%%
%% Key result: The simple spectral action V_eff produces an α²-scale
%% correction but NOT with coefficient 1/√8. The scale is confirmed;
%% the coefficient requires a separate mechanism.
%% ================================================================

\section{Appendix S: Numerical investigation of the second-order
  electromagnetic correction}
\label{app:B1computation}

This appendix reports a direct numerical computation of the
second-order electromagnetic variation
$V_{\mathrm{eff}}^{(2)}(R_0)$ of the effective potential at the
self-dual radius $R_0=\tfrac{1}{2}$, and compares it with the
coefficient $1/\sqrt{8}$ required by the SFST formula.  The result
resolves a partial sub-question of Barrier~B1 while making the
remaining gap more precise.

%% ----------------------------------------------------------------
\subsection{Setup and definitions}
\label{ssec:B1setup}

Expand the effective potential in powers of the electromagnetic
coupling:
\begin{equation}
  \label{eq:VeffExpand}
  V_{\mathrm{eff}}(R,\alpha)
    = V_{\mathrm{eff}}^{(0)}(R)
    + \alpha\,V_{\mathrm{eff}}^{(1)}(R)
    + \alpha^2\,V_{\mathrm{eff}}^{(2)}(R)
    + \mathcal{O}(\alpha^3).
\end{equation}
The first-order term $V_{\mathrm{eff}}^{(1)}(R_0)=0$ by the universal
tadpole cancellation (Appendix~R of the main text and Appendix~Q of
this supplement).  The present appendix concerns
$V_{\mathrm{eff}}^{(2)}(R_0)$.

\paragraph{Model operator.}
We use the simplified model of a phase twist $\tau=\theta/(2\pi)$ in
one torus direction, acting on the rank-$4$ spinor bundle with
charge-matrix weight $\operatorname{tr}(Q_{\mathrm{em}}^2)=\tfrac{4}{3}$
(the value for $V_C=\mathbf{3}\oplus\bar{\mathbf{3}}$ established in
Appendix~R).  The one-dimensional twisted heat kernel is
\[
  K_{1\mathrm{D}}(t,\tau,R)
    = \sum_{n\in\mathbb{Z}} e^{-t(n+\tau)^2/R^2},
\]
and the five-dimensional spectral action at twist $\tau$ is
\begin{equation}
  V_{\mathrm{eff}}(R,\tau)
    = \int_0^\infty \frac{dt}{t}\,
      \Bigl[K_{1\mathrm{D}}(t,\tau,R)^5
           + K_{\mathrm{wind}}(t,R,\alpha')^5\Bigr].
\end{equation}
The second-order coefficient in the $\tau$-expansion is
\begin{equation}
  \label{eq:Veff2def}
  V_{\mathrm{eff}}^{(2)}(R)
    = \frac{r\,\operatorname{tr}(Q_{\mathrm{em}}^2)}{(2\pi)^2}
      \int_0^\infty \frac{dt}{t}\,
      \left.\frac{\partial^2 K_{1\mathrm{D}}}{\partial\tau^2}
      \right|_{\tau=0}
      K_{1\mathrm{D}}(t,0,R)^4,
\end{equation}
where $r=4$ is the spinor rank, and the cross terms vanish because
$\partial_\tau K_{1\mathrm{D}}|_{\tau=0}=0$ by lattice antisymmetry.
The second derivative of the heat kernel is
\begin{equation}
  \left.\frac{\partial^2 K_{1\mathrm{D}}}{\partial\tau^2}
  \right|_{\tau=0}
    = \sum_{n\in\mathbb{Z}}
      \left(\frac{-2t}{R^2}+\frac{4t^2 n^2}{R^4}\right)
      e^{-tn^2/R^2}.
\end{equation}

%% ----------------------------------------------------------------
\subsection{Numerical results}
\label{ssec:B1results}

All integrals are evaluated with \texttt{mpmath} at 50-decimal-digit
precision, using log-spaced quadrature with 100 points over
$t\in[10^{-3},10^3]$ and lattice sums truncated at $|n|\leq 18$
(convergence verified).  The five-point finite-difference stencil
\[
  f''(x_0)
    \approx \frac{-f(x_0+2h)+16f(x_0+h)-30f(x_0)-16f(x_0-h)+f(x_0-2h)}{12h^2}
\]
with $h=10^{-3}$ is used for $\partial_R^2 V_{\mathrm{eff}}^{(0)}$.

\begin{table}[h]
\centering
\renewcommand{\arraystretch}{1.3}
\begin{tabular}{lrl}
\toprule
Quantity & Numerical value & Note \\
\midrule
$V_{\mathrm{eff}}^{(0)}(R_0)$ &
  $1.058\times 10^7$ & geometric baseline \\
$V_{\mathrm{eff}}^{(0)}{}''\!(R_0)$ &
  $5.189\times 10^8$ & $>0$, confirms global minimum \\
$V_{\mathrm{eff}}^{(2)}(R_0)$ &
  $-2.765\times 10^3$ & negative: EM lowers the action \\
$V_{\mathrm{eff}}^{(2)}(R_0)/V_{\mathrm{eff}}^{(0)}(R_0)$ &
  $-2.613\times 10^{-4}$ & relative second-order correction \\
$X\equiv V_{\mathrm{eff}}^{(2)}/(\alpha^2\,V_{\mathrm{eff}}^{(0)})$ &
  $-4.907$ & dimensionless coefficient \\
\midrule
$1/\sqrt{8}$ (SFST target) &
  $+0.3536$ & required coefficient \\
$X/(1/\sqrt{8})$ &
  $-13.9$ & \textbf{discrepancy} \\
\bottomrule
\end{tabular}
\caption{Numerical results for the second-order electromagnetic
  variation of the effective potential at $R_0=\tfrac{1}{2}$,
  $\alpha'=\tfrac{1}{4}$.  All computations in \texttt{mpmath} at
  50-digit precision.}
\label{tab:B1results}
\end{table}

%% ----------------------------------------------------------------
\subsection{Interpretation}
\label{ssec:B1interp}

\paragraph{What the computation establishes (Tier~1).}

\begin{enumerate}[label=(\roman*)]
  \item The second-order EM correction $V_{\mathrm{eff}}^{(2)}(R_0)$
        is nonzero: the quadratic spectral-action response to the
        electromagnetic coupling exists and is computable.

  \item The scale of the correction is of order $\alpha^2$:
        $|V_{\mathrm{eff}}^{(2)}/V_{\mathrm{eff}}^{(0)}|
        \approx 2.6\times10^{-4} \approx 5\,\alpha^2$.
        This confirms that the effective-potential approach produces
        corrections at the correct power of $\alpha$.

  \item The coefficient $X\approx -4.9$ differs from $1/\sqrt{8}$
        in both sign and magnitude.  The simple spectral-action
        model does \emph{not} generate the $1/\sqrt{8}$ coefficient.
\end{enumerate}

\paragraph{Why the sign is negative.}
The second derivative of $K_{1\mathrm{D}}$ contains a leading term
$-2t/R^2$, which integrates to a negative contribution.  The physical
interpretation is that electromagnetic fluctuations lower the
effective-action minimum, consistent with screening.  For the
observable mass-ratio correction to be positive (as required by the
formula), the matching map $\mathcal{M}[\ldots]$ must not be the
bare ratio $V_{\mathrm{eff}}^{(2)}/V_{\mathrm{eff}}^{(0)}$.  This is
consistent with the paper's treatment of $\mathcal{M}$ as an
unspecified but regulator-stable map.

\paragraph{What remains open (Tier~3).}
The factor $1/\sqrt{8}$ is not a geometric invariant of the
effective potential in the model studied here.  For it to emerge,
one of the following additional inputs is required:

\begin{enumerate}[label=(\alph*)]
  \item \textbf{Collective mode projection.}  The paper's §6
        motivates $1/\sqrt{8}$ as the normalised amplitude of a
        single isotropic collective mode from the eight-dimensional
        $\mathrm{SU}(3)$ adjoint sector.  This is a separate
        mechanism, distinct from the continuous spectral geometry of
        $V_{\mathrm{eff}}$.  The present computation neither confirms
        nor refutes it; it only shows that $V_{\mathrm{eff}}$ alone
        does not supply the factor.

  \item \textbf{Different observable.}  If the relevant quantity is
        a ratio of spectral determinants or zeta derivatives rather
        than $V_{\mathrm{eff}}^{(2)}/V_{\mathrm{eff}}^{(0)}$, the
        coefficient could differ.  A concrete alternative worth
        computing:
        $\delta\zeta'_{D^2}(0)/\zeta'_{D^2}(0)$ under the
        electromagnetic perturbation.

  \item \textbf{Full $Q_{\mathrm{em}}$-coupling in all directions.}
        The present model couples the EM twist to one torus direction.
        The full $Q_{\mathrm{em}}$ matrix couples to all five
        directions simultaneously.  A complete computation may alter
        the coefficient.
\end{enumerate}

\paragraph{Summary for B1.}
The computation confirms the $\alpha^2$ scale of the electromagnetic
correction.  It does not close Barrier~B1 (the $1/\sqrt{8}$
normalisation), but it sharpens the statement of the barrier: the
coefficient $1/\sqrt{8}$ is not a consequence of the leading
spectral-action geometry.  It must arise from one of the mechanisms
listed above, with~(a) being the candidate already proposed in the
main text.

\paragraph{Reproducibility.}
The computation uses: \texttt{mpmath} $\geq$ 1.3.0, 50-digit
precision, quadrature nodes 100, lattice truncation $|n|\leq 18$,
stencil step $h=10^{-3}$.  The key numerical values are
$V_{\mathrm{eff}}^{(2)}(R_0) = -2765.3$,
$V_{\mathrm{eff}}^{(0)}(R_0) = 1.058\times10^7$,
$V_{\mathrm{eff}}^{(0)}''\!(R_0) = 5.189\times10^8$.
A reviewer can reproduce these in fewer than two minutes of CPU time
using the scripts in the accompanying repository.
\end{document}
